\begingroup\scriptsize\setlength{\tabcolsep}{1.5pt}
\begin{tabular}{@{}p{2.7cm}rrrrrrrrl@{}}
\hline
Caso & $\tau$ (N m) & $Pe_r$ & $f$ (Hz) & $\Delta T$ (K) & $\eta_{\rm mecc}$ & $\eta_{\rm lim}$ & $t_{\rm spin}$ (s) & $F_{\rm scat}/F_g$ & Esito \\
\hline
Vaterite, sfera, acqua; $R=1\,\mu$m, $P=10$ mW & $2{,}5\times10^{-18}$ & 581{,}2 & 22{,}66 & 0{,}06 & $4{,}0\times10^{-14}$ & $7{,}6\times10^{-8}$ & -- & 299{,}5 & valido \\
Cu nero, sfera, acqua; $R=1{,}05\,\mu$m, $P=0{,}165$ mW & $4{,}4\times10^{-20}$ & 10{,}3 & 0{,}35 & 18{,}97 & $6{,}2\times10^{-16}$ & $6{,}1\times10^{-8}$ & -- & 1{,}3 & valido \\
Vaterite, sfera, vuoto $10^{-8}$; $R=1{,}05\,\mu$m, $P=0{,}316$ mW & $8{,}3\times10^{-20}$ & 19{,}5 & $1{,}0\times10^{12}$\textsuperscript{$\ast$} & 169{,}74 & $1{,}8\times10^{-3}$ & $7{,}7\times10^{-8}$ & 9{,}94 & 8{,}3 & solo formale \\
Vaterite, sfera, vuoto $10^{-8}$; $R=1{,}46\,\mu$m, $P=0{,}139$ mW & $5{,}0\times10^{-20}$ & 11{,}8 & $1{,}6\times10^{11}$\textsuperscript{$\ast$} & 61{,}87 & $2{,}8\times10^{-4}$ & $7{,}5\times10^{-8}$ & 63{,}04 & 1{,}3 & solo formale \\
\hline
\end{tabular}
\endgroup

\par\smallskip\noindent\footnotesize $^{\ast}$ In vuoto $f=\omega/(2\pi)$ con $\omega=\tau/\beta_{\rm gas}$ e' una frequenza formale, non una velocita' misurata; il rotore richiede una trappola. In acqua $t_{\rm spin}$ non e' applicabile (l'attrito domina) ed e' indicato con --. Le celle in acqua sono valide per i criteri, ma $\eta_{\rm lim}<10^{-7}$: non sono rotori fisici. Le celle in vuoto sono formali per lo stesso motivo. Qui $F_g=\rho Vg$ e $F_{\rm scat}=P/c$.
